\section{Empirical Evaluation}

We evaluate whether TreemapMix improves model Top-1 accuracy and calibration under
matched training recipes. The evaluation compares against mixed-sample
augmentation baselines.
CIFAR-100 and additional architectures are evaluated in
Appendices~\ref{app:cifar100-results} and~\ref{app:architecture-results}; optimizer
robustness is deferred to Appendix~\ref{app:optimizer-robustness}.

%\subsection{Experimental Setup}

We use a matched ImageNet-1K protocol to isolate the effect of
augmentation and supervision objective. Unless specified otherwise, each
experiment was executed on two NVIDIA H20 GPUs with a global batch size of 1024
and a per-GPU batch size of 512. Runs were executed on an H20 cluster whose
nodes each contained eight GPUs and 2 TB of system RAM; a custom scheduler
placed multiple independent two-GPU experiments on the same node to improve
utilization. Across the reported main and appendix experiments, total training
compute was approximately 5.5k GPU-hours. Peak observed host-memory usage
remained below 400 GB per node, indicating that the provisioned host memory
substantially exceeded the observed requirements; minimum host-memory
requirements were not separately profiled. ImageNet-1K is used as the primary
reference benchmark~\citep{dengImageNetLargescaleHierarchical2009}.

Our primary comparison uses four compact ViT/EVA02-family variants implemented
in \texttt{timm}~\citep{fangEVA02VisualRepresentation2024}. Unless specified
otherwise, each experiment is repeated with seeds 42, 43, and 44. For the balanced-data experiments, we use a balanced round-robin dataloader
that provides class-diverse source images for each composed sample. The
ImageNet-LT experiment instead preserves the natural class frequencies.
Appendix~\ref{app:appendix-aug-configs} gives an overview of the training configurations and hyperparameter settings, and Appendix~\ref{app:loss-code} provides reference implementations of the loss functions.

\paragraph{Comparison to Image Augmentation Baselines:}

We first test whether TreemapMix improves ImageNet-1K performance under the
matched 110-epoch protocol. For the main comparison, we use $K=4$ for
TreemapMix-SCE and $K=6$ for TreemapMix-PL, together with a cosine schedule
for $\alpha$ from $0.1$ to $0.5$. These values were selected from preliminary
tuning runs and then kept fixed across all backbones in the main comparison.
For TreemapMix-mixed, we set $\beta=0.5$, giving equal weight to the
TreemapMix-SCE and TreemapMix-PL losses, and use the same $\alpha$ schedule
as the other TreemapMix variants at $K=4$.

For TreemapMix, we evaluate three supervision variants: area-proportional soft cross-entropy (TreemapMix-SCE), the Plackett-Luce ranking loss (TreemapMix-PL), and their convex combination (TreemapMix-mixed). We compare these variants against Mixup, CutMix, Mixup+CutMix, RICAP, FMix, GridMix, ResizeMix, SaliencyMix, SmoothMix, TokenMix, and TLA.
The seven additional methods, FMix through TLA, are reimplementations of
methods available in OpenMixup~\citep{liOpenMixupOpenMixup2024}. For each,
we evaluated eight hyperparameter configurations on ViT-Wee, selected the
best-performing configuration, and fixed it across the four backbones.
Figure~\ref{fig:in1k-main-results-top1} summarizes the ImageNet-1K
results for Top-1 accuracy and ECE using 15 bins; higher is better for Top-1 accuracy,
and lower is better for ECE. Appendix
Appendix~\ref{app:in1k-confidence-results} reports the full Top-1 accuracy
and ECE tables for the main ImageNet-1K comparison.
\begin{figure*}[t]
  \centering
  \includegraphics[width=0.9\textwidth]{../evaluation/data/processed/figures/per_experiment/in1k_per_experiment_short_legend.pdf}\\[0pt]
  \includegraphics[width=0.9\textwidth]{../evaluation/data/processed/figures/per_experiment/in1k_per_experiment_short.pdf}
\caption{ImageNet-1K main comparison under the matched 110-epoch protocol.
Across backbones, TreemapMix-PL gives the strongest Top-1 accuracy, while
TreemapMix-SCE gives the lowest ECE, showing a consistent accuracy-calibration
trade-off between the two objectives. Bars show the mean over three seeds, and
error bars denote $\pm 1$ standard deviation. ECE is computed with 15 bins and
reported in percentage points.}
  \label{fig:in1k-main-results-top1}
\end{figure*}

Across all four backbones, TreemapMix-PL gives the strongest Top-1 accuracy, while TreemapMix-SCE is the most calibration-oriented variant. The mixed objective usually falls between them, suggesting that the two supervision signals are not simply additive. Among the non-TreemapMix baselines, the strongest method depends on model size: single-image augmentation is strongest on ViT-Wee, SaliencyMix is strongest on ViT-Little, and ResizeMix is strongest on ViT-Medium and ViT-Betwixt.

Calibration metrics show the opposite preference from Top-1 accuracy.
TreemapMix-SCE gives the lowest ECE, supporting an accuracy-calibration
trade-off between rank-oriented and probability-matching supervision.
One explanation is that PL mainly constrains relative logit order, whereas SCE also penalizes probability magnitudes. This can improve discrimination while leaving calibration errors less constrained.

\paragraph{Long-Horizon Training Behavior:}
\label{sec:long-horizon}
We test whether TreemapMix remains useful under longer training with
ViT-Wee and ViT-Little. Table~\ref{tab:in1k-long-horizon-results} compares
single-image augmentation, Mixup+CutMix, TreemapMix-SCE, and TreemapMix-PL
using seeds 42, 43, and 44. The schedule comprises 320 total epochs:
20 warmup epochs followed by 300 regular epochs.

TreemapMix-PL gives the highest mean Top-1 accuracy, while TreemapMix-SCE
gives the lowest mean ECE on both backbones. Relative to Mixup+CutMix,
PL improves accuracy by 0.91 percentage points on ViT-Wee and 0.21 points
on ViT-Little, smaller gains than under the 110-epoch schedule. The PL
advantage over single-image augmentation is 0.27 and 0.83 points,
respectively. These runs reuse the TreemapMix settings selected for the
shorter schedule and support persistence of the gains rather than a fully
tuned long-horizon comparison.

\begin{table*}[t]
  \caption{ImageNet-1K long-horizon results for ViT-Wee and ViT-Little under a 320-epoch schedule. Accuracy is reported in percent, and ECE in percentage points using 15 bins. Values are mean $\pm$ one sample standard deviation over seeds 42, 43, and 44. Bold marks the best mean for each model and metric.}
  \label{tab:in1k-long-horizon-results}
  \centering
  \footnotesize
  \setlength{\tabcolsep}{4pt}
  \begin{tabular}{@{}lcccc@{}}
    \toprule
     & \multicolumn{2}{c}{ViT-Wee} & \multicolumn{2}{c}{ViT-Little} \\
    \cmidrule(lr){2-3} \cmidrule(lr){4-5}
    Method & Acc. $\uparrow$ & ECE $\downarrow$ & Acc. $\uparrow$ & ECE $\downarrow$ \\
    \midrule
    Single-Image Aug & $80.37 \pm 0.08$ & $3.89 \pm 0.16$ & $81.12 \pm 0.04$ & $3.39 \pm 0.06$ \\
    Mixup+CutMix & $79.73 \pm 0.15$ & $6.87 \pm 0.64$ & $81.74 \pm 0.09$ & $5.47 \pm 0.52$ \\
    TreemapMix-SCE & $80.41 \pm 0.04$ & $\mathbf{3.23 \pm 0.49}$ & $81.58 \pm 0.33$ & $\mathbf{2.77 \pm 1.42}$ \\
    TreemapMix-PL & $\mathbf{80.64 \pm 0.16}$ & $6.94 \pm 0.14$ & $\mathbf{81.95 \pm 0.34}$ & $7.44 \pm 0.54$ \\
    \bottomrule
  \end{tabular}
\end{table*}


\paragraph{Computational Performance:}
\label{sec:runtime}
Data-loader throughput and host-memory measurements are reported in
Appendix~\ref{app:runtime-overhead}.

\paragraph{Additional Datasets and Architectures:}
The appendix reports long-tailed recognition results on ImageNet-LT
(Appendix~\ref{app:long-tail-results}), results on CIFAR-100
(Appendix~\ref{app:cifar100-results}), and results for ConvNeXtV2-Nano,
DeiT-Tiny, and Swin-Tiny on ImageNet-1K
(Appendix~\ref{app:architecture-results}).
TreemapMix-SCE achieves the lowest ECE on all these backbones.
TreemapMix-PL achieves the highest Top-1 accuracy on all three additional
ImageNet-1K architectures and three of the four CIFAR-100 backbones;
ResizeMix leads on CIFAR-100 with ViT-Little.
